\documentclass[10pt,reqno]{amsart}

\usepackage[T1]{fontenc}
\usepackage[utf8]{inputenc}
\usepackage{newtxtext,newtxmath}
\usepackage{microtype}
\usepackage{mathtools}
\usepackage{enumitem}
\usepackage[nameinlink,noabbrev]{cleveref}
\usepackage[a4paper,margin=28mm,headheight=15pt]{geometry}

\setlength{\parskip}{2pt}
\setlist[enumerate]{leftmargin=7mm,itemsep=2pt,topsep=3pt}
\emergencystretch=2em

\newtheorem{theorem}{Theorem}[section]
\newtheorem{proposition}[theorem]{Proposition}
\newtheorem{lemma}[theorem]{Lemma}
\newtheorem{corollary}[theorem]{Corollary}
\theoremstyle{definition}
\newtheorem{remark}[theorem]{Remark}

\title[A symmetric obstruction to the pseudo-triangulation road map]
{A Half-Turn Symmetric Octagon Obstructing the Standard Pseudo-Triangulation Road Map}

\author{Argus Open Research Project}
\address{Argus Open Research Project}
\date{August 31, 2026}

\subjclass[2020]{Primary 52C25; Secondary 68U05, 05C10}
\keywords{pseudo-triangulation, carpenter's rule, expansive motion, symmetry, rigidity}

\begin{document}

\begin{abstract}
An official AIM workshop summary already records a negative answer to the
symmetry question for the pseudo-triangulation algorithm, but gives no
explicit witness.  We give a rational simple nonconvex octagon for which every legal first step
of the standard pointed-pseudo-triangulation road map breaks the original
labelled half-turn symmetry.  Every convex-hull edge belongs to a half-turn
orbit of size two.  The road map deletes exactly one hull edge from a pointed
pseudo-triangulation, leaving its symmetric partner as a rigid bar.  Any
half-turn-preserving motion must keep the two paired distances equal, so the
deleted distance is constant as well.  All edge lengths of the original
infinitesimally rigid completion are then fixed, and no nontrivial first
motion is possible.  This does not obstruct the distinct canonical
symmetry-preserving expansive motion of Connelly, Demaine, and Rote, nor a
modified pseudo-triangulation method that deletes an entire symmetry orbit.
\end{abstract}

\maketitle

\section{Question and status}

AIM Problem~13 asks whether the pseudo-triangulation algorithm for the
carpenter's-rule problem gives an unfolding that preserves the symmetries of
the initial linkage.  We use ``the pseudo-triangulation algorithm'' in its
published standard sense: complete the polygon to a pointed
pseudo-triangulation, delete one eligible convex-hull edge, and follow the
resulting one-degree-of-freedom expansive mechanism.

The AIM workshop summary, publicly archived by September~18, 2015, states
that this question is false and separately notes that the construction of
Connelly, Demaine, and Rote does preserve symmetry.  It supplies neither a
witness nor a proof.  Accordingly, the negative status is not new: the
contribution of this note is an explicit rational witness and a
self-contained obstruction argument for the standard road map.  No priority
claim is made for the negative answer itself.

\section{The symmetric octagon}

Let \(P=(p_0,\ldots,p_7)\) be the polygon with vertices, in cyclic order,
\[
\begin{aligned}
p_0&=(2,1),             &p_1&=(0,\tfrac12),\\
p_2&=(-2,1),            &p_3&=(-1,0),\\
p_4&=(-2,-1),           &p_5&=(0,-\tfrac12),\\
p_6&=(2,-1),            &p_7&=(1,0).
\end{aligned}
\]
The half-turn
\[
 h(p_i)=p_{i+4}=-p_i
\]
is a labelled symmetry of the polygonal linkage.

\begin{lemma}\label{lem:geometry}
The polygon \(P\) is simple and nonconvex.  Its convex hull has vertices
\[
 p_0,p_2,p_4,p_6,
\]
and none of its four hull edges is a polygon edge.  The hull-edge orbits
under \(h\) are
\[
 p_0p_2\longleftrightarrow p_4p_6,
 \qquad
 p_2p_4\longleftrightarrow p_6p_0.
\]
\end{lemma}

\begin{proof}
Direct determinant calculations show that the signed turns alternate
between \(-2\) and \(3/2\).  Thus the polygon is nonconvex and has no
collinear consecutive vertices.  Checking the twenty pairs of nonincident
edges shows that no pair intersects.  The four extreme vertices are exactly
\((\pm2,\pm1)\), while the remaining four vertices lie strictly inside their
convex hull.  The displayed half-turn sends each hull edge to the opposite
one, and no hull edge joins consecutive polygon vertices.
\end{proof}

\section{The standard road-map step}

A pointed pseudo-triangulation on \(n\) vertices is a maximal planar pointed
graph and has \(2n-3\) edges.  A simple polygon, regarded as a noncrossing
pointed edge set, can be completed to such a graph.  The standard road-map
convexification step chooses a pointed-pseudo-triangulation completion
\(T\), deletes one convex-hull edge that is not a polygon edge, pins a
remaining edge, and follows the resulting one-degree-of-freedom expansive
mechanism until the next alignment event.

For the octagon in \Cref{lem:geometry}, every completion \(T\) contains the
eight polygon edges and all four convex-hull edges.  These twelve edges are
distinct.  Since \(T\) has
\[
 2\cdot8-3=13
\]
edges, the completion contains only one additional edge.  The argument below
does not depend on that edge or on the choice of completion.

\section{The missing-edge obstruction}

\begin{theorem}\label{thm:main}
No legal first step of the standard pointed-pseudo-triangulation road map
for \(P\) preserves its labelled half-turn symmetry.
\end{theorem}

\begin{proof}
Fix an arbitrary pointed-pseudo-triangulation completion \(T\), and let
\(e=p_ap_b\) be the hull edge deleted in the first road-map step.  By
\Cref{lem:geometry},
\[
 h(e)=p_{a+4}p_{b+4}\neq e.
\]
The partner \(h(e)\) is another mandatory hull edge of \(T\), and it remains
a rigid bar in the mechanism \(T-e\).

Suppose that a motion \(q_i(t)\) of \(T-e\) preserved the original labelled
half-turn.  We allow both its center and the whole framework to move.  Thus
there would be a point \(c(t)\) such that
\[
 q_{i+4}(t)=2c(t)-q_i(t)
 \qquad(i\in\mathbb Z/8\mathbb Z).
\]
It follows that
\[
 q_{a+4}(t)-q_{b+4}(t)
 =-\bigl(q_a(t)-q_b(t)\bigr),
\]
and hence
\[
 \lVert q_{a+4}(t)-q_{b+4}(t)\rVert
 =\lVert q_a(t)-q_b(t)\rVert.
\]
The left-hand side is the length of the retained bar \(h(e)\), so it is
constant.  The missing distance \(e\) is therefore constant as well.
Consequently every edge length of the full pointed pseudo-triangulation
\(T\) is constant during the motion.

Pointed pseudo-triangulations are infinitesimally rigid and hence locally
rigid.  The congruence orbit of the initial realization is therefore both
closed and locally open inside any continuous path preserving all edge
lengths of \(T\).  Since the parameter interval is connected, the whole path
remains in that congruence orbit.  After the road map pins an edge, even this
ambient rigid-body motion is removed.

The first road-map segment is required to be the nontrivial expansive motion
of the one-degree-of-freedom mechanism \(T-e\).  It therefore cannot preserve
the half-turn.  Since \(T\) and \(e\) were arbitrary, every legal first step
breaks the symmetry.
\end{proof}

\begin{remark}
The infinitesimal form is even shorter.  Orient the unique mechanism flex so
that the deleted hull distance has positive derivative.  Half-turn symmetry
would make this derivative equal to that of the retained partner bar, which
is zero.
\end{remark}

\section{Scope}

\begin{remark}\label{rem:scope}
The theorem concerns the standard road map in which exactly one hull edge is
deleted at each step.  It does not rule out a modified method that deletes an
entire symmetry orbit and analyzes a higher-dimensional mechanism.  It also
does not assert that no symmetric pointed pseudo-triangulation exists.
\end{remark}

There is no contradiction with the symmetry-preserving carpenter's-rule
motion of Connelly, Demaine, and Rote.  Their construction selects the unique
solution of a global optimization problem; applying a symmetry to the
solution gives the same solution by uniqueness.  That canonical motion is a
different mechanism from the finite sequence of one-edge-deletion
pseudo-triangulation motions considered here.

\appendix

\section{Checklist for verification}

\begin{enumerate}
  \item Check polygon simplicity and the alternating signed turns.
  \item Recompute the four hull vertices and verify that no hull edge is a
  polygon edge.
  \item Verify the two size-two half-turn orbits of hull edges.
  \item Check that every pointed-pseudo-triangulation completion contains all
  polygon and hull edges and has thirteen edges.
  \item For each possible deleted hull edge, verify that its half-turn partner
  remains a rigid bar.
  \item Use the moving-center half-turn equation to prove equality of the two
  paired distances.
  \item Apply local rigidity of the full completion and keep the standard
  road map separate from other expansive motions.
\end{enumerate}

\begin{thebibliography}{9}

\bibitem{CDR}
R.~Connelly, E.~D. Demaine, and G.~Rote,
\emph{Straightening polygonal arcs and convexifying polygonal cycles},
Discrete Comput. Geom. \textbf{30} (2003), 205--239.

\bibitem{Streinu}
I.~Streinu,
\emph{Pseudo-triangulations, rigidity and motion planning},
Discrete Comput. Geom. \textbf{34} (2005), 587--635.

\bibitem{AIM}
American Institute of Mathematics,
\emph{Configuration spaces of linkages},
workshop problem list and summary, October 13--17, 2014, Problem~13.

\end{thebibliography}

\end{document}
